\section{Results}\label{sec:results}

The results answer the three questions of Sect.~\ref{sec:intro} in turn:
\begin{itemize}[leftmargin=*,nosep]
  \item \textbf{Q1, perfect model:} Sects.~\ref{sec:q1}--\ref{sec:q1-density};
  \item \textbf{Q2, model error:} Sects.~\ref{sec:q2}--\ref{sec:q2-represent};
  \item \textbf{Q3, form of the prior:} Sects.~\ref{sec:q3} and~\ref{sec:composition}.
\end{itemize}
The transfer to the QG ocean and the value of true-system data close the section. Unless stated, numbers are Lorenz-96 RMSE on the regular observing system.

\subsection{Overview}\label{sec:overview}
Table~\ref{tab:overview} gives the RMSE of every scheme in both cases and on both
observing systems. With a perfect model, the learned schemes lead by 30 to
45\,\%:
\begin{itemize}[leftmargin=*,nosep]
  \item DirectUNet 0.340 and PredictStateCFM 0.341 on the regular grid;
  \item the best DA scheme, the ETKS, 0.497;
  \item Strong-4DVar, which assimilates with the true equations, 0.703.
\end{itemize}
The following subsections ask where the DA deficit comes from, with a perfect model and then under model
error, and whether the form of the prior matters.

\begin{table}[t]
\centering\small
\caption{RMSE on the 24 observed channels, 200 test windows; learned rows: 1200 epochs, 3 seeds pooled per window. The italic row is a sensitivity row. Sources: \texttt{reports/l96/outputs/p1\_l96\_benchmark.md}, \texttt{docs/results/l96\_weak4dvar.md}.}
\label{tab:overview}
\begin{tabular}{lcccc}
\toprule
 & \multicolumn{2}{c}{\Sz{} (perfect model)} & \multicolumn{2}{c}{\So{} (model error)} \\
\cmidrule(lr){2-3}\cmidrule(lr){4-5}
scheme & regular & random & regular & random \\
\midrule
ETKF & 0.610 & 0.679 & 1.409 & 1.407 \\
EnKF & 0.641 & 0.706 & 1.416 & 1.484 \\
ETKS & 0.497 & 0.572 & 1.338 & 1.347 \\
Strong-4DVar & 0.703 & 0.742 & 1.436 & 1.444 \\
Weak-4DVar & 0.695 & \needsrun{} & 1.064 & \needsrun{} \\
\emph{ETKS, 100 members} & \emph{0.393} & \emph{0.463} & \emph{1.338} & \emph{1.335} \\
\midrule
DirectUNet & 0.340 & 0.450 & 0.339 & 0.444 \\
PredictStateCFM & 0.341 & 0.443 & 0.338 & 0.447 \\
VanillaCFM & 0.351 & 0.462 & 0.349 & 0.468 \\
SDA1 & 0.464 & 0.572 & 0.462 & 0.585 \\
SDA2 & 0.453 & 0.550 & 0.465 & 0.575 \\
SDA3-fix & 0.455 & 0.562 & 0.455 & 0.573 \\
hybrid DirectUNet $\to$ SDA3-fix & \textbf{0.293} & \textbf{0.367} & \textbf{0.289} & \textbf{0.367} \\
\bottomrule
\end{tabular}
\end{table}

\subsection{Error budget with a perfect model}\label{sec:q1}\label{sec:q1-budget}
\paragraph{Irreducible.}
The best scheme, the hybrid, reaches 0.293 on the regular grid: the irreducible
error is at most this. We have no lower bound, so every share below is a share of
the excess over an upper bound, not over the true $\mathrm{mmse}$.

\paragraph{Restriction: filtering.}
The ETKS removes 16 to 19\,\% of the ETKF's RMSE (0.610 $\to$ 0.497 regular,
0.679 $\to$ 0.572 random). To check that this gain is a restriction, we run the
filter and the smoother on the same ensemble realisation and evaluate both sides
of the identity
$\mathrm{MSE}_F - \mathrm{MSE}_S = \mathrm{var}_F - \mathrm{var}_S = \E\|m_F - m_S\|^2$.
It holds approximately with 100 members --- an MSE gap of 0.162 against a
variance gap of 0.176 and a mean shift of 0.233 --- and less well with 30
(0.159, 0.227, 0.296). The deviation shrinks as the ensemble grows, as it should
if it is an approximation error of the ensemble.

\paragraph{Approximation: finite ensemble.}
With 100 members and a re-selected inflation, the ETKS goes from 0.497 to 0.393
(regular) and from 0.572 to 0.463 (random): 38\,\% of its mean-squared error was
finite-ensemble approximation. Against the best learned scheme, the DA deficit
falls from $1.46\times$ to $1.16\times$.

\paragraph{Approximation: optimisation and training.}
\begin{itemize}[leftmargin=*,nosep]
  \item \textbf{4D-Var.} At the benchmark's 10 L-BFGS iterations per sub-window,
    the hard-constraint 4D-Var is far from its own optimum: 0.721 against 0.545
    on the validation windows with 40 iterations.
  \item \textbf{DirectUNet.} Training it for 1200 instead of 400 epochs takes it
    from 0.380 to 0.340, and three times more training windows add only 0.006
    (0.334). The learned approximation term is close to saturation at our
    budget.
\end{itemize}

\paragraph{Restriction: context.}
The unconditional SDA prior (SDA1, 0.464) ignores $\thv$. Conditioning it on
$\thv$ gains 2\,\% (SDA2 0.453, SDA3-fix 0.455). \needsrun{The forcing $\zv$ as a
conditioning input of SDA, and $(\thv, \zv)$ as inputs of DirectUNet, complete
this row.}

\paragraph{Structural misspecification.}
What remains of the DA excess after these contrasts is structural: the mode
instead of the mean, a hard constraint and Gaussian updates. It is largest for
Strong-4DVar.
\todo{Table~5: the budget in pooled MSE per scheme class, both observing
systems, with the contrast named in each cell.}

\paragraph{Reading.}
With a perfect model, the deficit of classical DA is mostly \emph{honest}. It is
made of a restriction (filtering) and an approximation (30 members, a short
optimisation), and it shrinks with smoothing, with members and with iterations.
None of it requires a better model.

\subsection{Reported uncertainty with a perfect model}\label{sec:q1-uncertainty}
Table~\ref{tab:fms-s0} gives $\FMS_\tau$ for the main schemes at \Sz{} on the regular grid.
\begin{itemize}[leftmargin=*]
  \item \textbf{The smoother buys its mean with an over-confident ensemble.} It
    beats the filter at $\tau = 0$ (0.129 vs 0.187) and still at $\tau = 0.5$,
    but loses at $\tau = 0.75$ (0.059 vs 0.050). Its calibration loss there is
    34\,\% against 14\,\%.
  \item \textbf{The flows are nearly calibrated} (calibration loss 0--3\,\%,
    pooled spread/skill 0.84--0.94).
  \item \textbf{The SDA samplers are under-dispersed} (7--18\,\%), as expected
    from their tempered likelihood.
\end{itemize}
\needsrun{The filter/smoother identity along $\tau$ (score of the ETKS's own
filter pass) and the 100-member rows of this table.}

\begin{table}[t]
\centering\small
\caption{$\FMS_\tau$ (standardised units, Gaussian form) and calibration loss at $\tau = 0.75$, \Sz, regular grid. Source: \texttt{docs/results/l96\_p1\_fm\_score.md}.}
\label{tab:fms-s0}
\begin{tabular}{lcccc}
\toprule
scheme & $\tau = 0$ & $\tau = 0.5$ & $\tau = 0.75$ & calibration loss \\
\midrule
ETKF & 0.187 & 0.125 & 0.050 & 14\,\% \\
ETKS & 0.129 & 0.107 & 0.059 & 34\,\% \\
PredictStateCFM & 0.054 & 0.040 & \textbf{0.0195} & 2\,\% \\
VanillaCFM & 0.057 & 0.042 & 0.0196 & 0\,\% \\
SDA3-fix & 0.088 & 0.067 & 0.034 & 14\,\% \\
hybrid & \textbf{0.039} & \textbf{0.034} & 0.0200 & 16\,\% \\
\bottomrule
\end{tabular}
\end{table}

\figplaceholder{Fig.~3}{$\FMS_\tau$ curves per scheme, \Sz, regular and random observing systems, with window-bootstrap intervals.}

\subsection{Dependence on the observing system}\label{sec:q1-density}
Figure~4 shows how the schemes depend on the number of observations.
\begin{itemize}[leftmargin=*]
  \item \textbf{Crossover.} DA is best when a window has at most about ten
    observation times; the learned schemes win from about twenty.
  \item \textbf{Transfer between observing systems.} The ratio of random to
    regular RMSE is 1.06 for Strong-4DVar, 1.15 for the ETKS, 1.21--1.24 for
    SDA and 1.30--1.32 for the CFMs and DirectUNet. Keeping $\Hop$ and $\mathbf{R}$
    explicit buys robustness to the observing system.
  \item \textbf{Beyond the training range} (from 300 to 1000 observations),
    DirectUNet degrades from 0.17 to 0.34, while SDA, whose likelihood is
    explicit, keeps improving from 0.30 to 0.25.
\end{itemize}
In the language of the chain, an amortised map is misspecified once the
observing system leaves its training distribution; a factorised scheme is not.

\figplaceholder{Fig.~4}{RMSE vs number of observation times and of observed fast channels, \Sz, all schemes.}

\subsection{Error budget under model error}\label{sec:q2}
Under model error the DA errors grow $2.0$--$2.7\times$ in RMSE: ETKF 1.409, EnKF
1.416, ETKS 1.338, Strong-4DVar 1.436. For the ETKS the mean-squared error is about seven
times its \Sz{} value. The contrasts that explained the \Sz{} deficit no longer
help.
\begin{itemize}[leftmargin=*]
  \item \textbf{Ensemble size.} 100 members change nothing (1.338 $\to$ 1.338).
  \item \textbf{Smoothing.} It gains only 5\,\%.
  \item \textbf{The filter/smoother identity fails.} The MSE gap is 0.237, against
    a variance gap of 0.534 and a mean shift of 0.913.
\end{itemize}
The excess is misspecification. The learned rows are flat (DirectUNet 0.339,
hybrid 0.289), but they are an oracle bound: they never saw the biased model
(Sect.~\ref{sec:design}).

\subsection{Reported uncertainty under model error}\label{sec:q2-uncertainty}
The large-$\tau$ score shows misspecification even where the RMSE does not
single it out (Fig.~5).
\begin{itemize}[leftmargin=*]
  \item \textbf{DA calibration loss grows from \Sz{} to \So.} At $\tau = 0.75$
    it goes from 34 to 63\,\% for the ETKS, with an optimal variance rescaling
    $c^\star \ge 8$. For the ETKF it goes from 14 to 23\,\% on the regular grid
    and from 3 to 14\,\% on the random one.
  \item \textbf{The DA ensembles become short-tailed.} A kernel density of the
    members scores 17 to 30\,\% worse than a Gaussian of the same variance:
    the truth falls outside the ensemble more often than its spread allows.
  \item \textbf{The learned rows are invariant}, with one exception: SDA2,
    trained with clean parameters and conditioned on the biased ones, loses less
    at \So{} (7 $\to$ 3\,\%) because the bias accidentally widens its samples.
  \item \textbf{Aggregate calibration is not calibration.} On the random system
    at \So, the ETKF's pooled spread/skill is 0.99 --- nominally calibrated ---
    yet it loses 12\,\% at $\tau = 0.75$. A scalar inflation can match the total
    error, but not where the error occurs.
\end{itemize}

\figplaceholder{Fig.~5}{$\FMS_\tau$ at \Sz{} and \So{} per scheme; calibration loss vs $\tau$.}

\subsection{Marginal value of observations}\label{sec:q2-mvo}
The clearest symptom is what additional observations buy (Fig.~6).
\begin{itemize}[leftmargin=*]
  \item \textbf{Strong-4DVar.} Going from 15 to 30 observations per window, its
    RMSE gain collapses $7.0\times$ under model error. It then stays on a floor
    set by the model: 1.43 at 30 observations, 1.38 at 1000.
  \item \textbf{Filters.} They keep most of the value of observations: under
    model error their gain from the same doubling shrinks only $1.8$--$1.9\times$,
    against $7.0\times$ for Strong-4DVar.
  \item \needsrun{Weak-4DVar and the learned schemes on the same curve.}
\end{itemize}
A hard constraint turns model error into a misspecification that no amount of
data repairs.

\figplaceholder{Fig.~6}{RMSE vs number of observation times, \Sz{} and \So, per scheme.}

\subsection{Representing model uncertainty}\label{sec:q2-represent}
If misspecification is the problem, representing the uncertainty of the model
should convert part of it back into a restriction.
\begin{itemize}[leftmargin=*]
  \item \textbf{Weak-4DVar.} We tune its model-error scale per case on the
    validation windows. At \Sz{} it gains nothing over the hard constraint
    (0.695 vs 0.703), as it should with a perfect model. At \So{} it removes
    26\,\% of the error (1.064 vs 1.436) and becomes the best DA scheme, ahead of
    the ETKS (1.338). It still remains three times the learned bound: an
    additive per-step error with a scalar scale cannot absorb a parametric bias.
  \item \textbf{Learned priors.} SDA2, trained with clean parameters, leaks the
    \So{} bias into its prior, and the hybrid built on it degrades
    (0.310 $\to$ 0.333 with 400-epoch priors). SDA3-fix, trained with noisy
    parameters, stays flat (0.312 $\to$ 0.307).
  \item \needsrun{A perturbed-parameter ETKF, and the $\delta$ sweep: the
    conversion should succeed while the bias stays within the stated
    uncertainty ($\delta \le \sigma$) and fail beyond.}
\end{itemize}

\figplaceholder{Fig.~7}{RMSE and calibration loss vs bias $\delta$, per model-uncertainty arm.}

\subsection{Model-derived and learned priors}\label{sec:q3}
The posterior~\eqref{eq:posterior} is a likelihood times a prior. Model-based DA
derives the prior from the equations; SDA learns it from data and keeps the same
explicit likelihood. Comparing them isolates the form of the prior.

\paragraph{Same likelihood, two priors.}
At \Sz{} the learned joint-window prior (SDA, 0.453--0.464 on the regular grid)
beats the true-model hard constraint (Strong-4DVar, 0.703) and the 30-member
smoother (ETKS, 0.497). It does not beat the 100-member smoother (0.393). The
ordering is therefore set by the approximation term of the model-based scheme,
not by the fidelity of its prior --- which, at \Sz, is exact. Along $\tau$, the
learned prior also gives the better-calibrated posterior: calibration loss 7 to
18\,\% for SDA against 34\,\% for the 30-member smoother (Table~\ref{tab:fms-s0}).

\figplaceholder{Fig.~8}{Model-derived vs learned priors at matched likelihood: RMSE and $\FMS_\tau$, \Sz{} and \So.}

\paragraph{A better prior is not a better estimate.}
Two observations make the point.
\begin{itemize}[leftmargin=*]
  \item Training the SDA prior three times longer (400 to 1200 epochs) improves the
    estimate by 7 to 9\,\%.
  \item The true equations, used as a hard prior, give the worst DA scheme at
    \Sz.
\end{itemize}
What matters is how much of the prior the scheme can actually use within its
approximation budget. This reading also resolves the controlled result that
motivated us~\cite{fablet4dvarnet}: a learned prior can beat the true model
because of the approximation and restriction terms of the scheme that uses the
true model, not because the learned prior is better.

\paragraph{Conditioning the prior.}
Conditioning the learned prior on the parameters is a restriction term at \Sz{}:
SDA1 0.464 vs SDA2 0.453 and SDA3-fix 0.455. Under model error it becomes a
misspecification channel for SDA2, which was trained with clean parameters and
degrades by 3\,\%, but not for SDA3-fix, trained with noisy ones.
\needsrun{Conditioning on the forcing $\zv$.}

\subsection{Composing a mean and an uncertainty}\label{sec:composition}
In the operator language of Sect.~\ref{sec:fms}, the hybrid composes DirectUNet's
mean with the anomaly of SDA3-fix's guided sampler. It has the best mean
everywhere (0.293 regular, 0.367 random, flat at \So) but not the best
distribution: the flows overtake it at $\tau = 0.75$ (0.0195 vs 0.0200 on the
regular grid, 0.025 vs 0.029 on the random one, standardised units). The score
separates the two components the hybrid composes: an amortised mean, accurate in
distribution, and an explicit-likelihood sampler whose tempered guidance makes it
over-confident.

\subsection{Transfer to the QG ocean}\label{sec:q1-qg}
\paragraph{Perfect model.}
On the QG case study (score = mean explained variance of the two
streamfunction and two potential-vorticity layers, higher is better; 100 test
windows of 30 days, forced dataset) the perfect-model ordering of the DA schemes
is the same as on Lorenz-96. Smoothing gains over filtering (EnKS 0.812 against
ETKF 0.774 and EnKF 0.761), and the variational scheme, which uses the whole
window, is best (strong-constraint 4D-Var 0.840). With a perfect model the
restriction of filtering is again the main reducible term among the DA schemes.
\needsrun{QG learned schemes (round 1) for the cross-class comparison.}

\paragraph{Model error.}
Under its realistic model error the QG case study repeats the Lorenz-96 pattern.
\begin{itemize}[leftmargin=*]
  \item \textbf{Hard constraint.} Strong-constraint 4D-Var, best at \Sz{}
    (0.840), becomes the worst scheme at \So{} with its \Sz{} settings (0.516).
  \item \textbf{Ensembles.} The ensemble schemes keep more (ETKF 0.579, EnKF
    0.596, EnKS 0.612).
  \item \textbf{Representing model uncertainty.} Re-tuning for \So{} --- a weak
    constraint for 4D-Var, an inflated observation error for the ensembles ---
    recovers part of the loss: 4D-Var 0.612, ETKF 0.633, EnKF 0.634, EnKS 0.665.
\end{itemize}
The smoother is the most robust scheme under model error on both systems, and
the hard constraint the least. Misspecification turns the best perfect-model
scheme into the worst.

\subsection{Value of true-system data}\label{sec:q2-oracle}
\needsrun{Learned schemes trained on simulations of the biased model, with
$\sigma = 0$ and 20\,\%; their gap to the learned-on-truth rows measures what
training on reanalyses or true-system data implicitly spends.}
